\section{Conclusion, Limitations, and Broader Impacts}
\label{limitation}
In this work, we propose an in-context RL method HDM that achieves both high effectiveness and efficiency in long-term tasks. The core idea of DM-H is to use sub-goal settings to leverage Mamba's ability to efficiently recall long sequence contexts while using the transformer to perform high-quality predictions. Unlike current in-context RL methods limited to short-term tasks, DM-H is also good at standard RL benchmarks, which typically have long-term sequences under the across-episodic context setting. On the Grid World, D4RL, and Tmaze benchmarks, we demonstrate that DM-H can outperform baselines in both efficiency and effectiveness.

Regarding limitations, our method has an important hyperparameter, which is the context length of the transformer "$c$." As "$c$" increases, Mamba model can scan multiple episodes with smaller context sizes, significantly improving the computational efficiency. On the contrary, as "$c$" decreases, Mamba can generate high-quality sub-goals from the context that is closer to the original episode. However, this reduces the short-term context available to the transformer when making predictions. Therefore, a meaningful future direction is for "$c$" to adapt autonomously to different tasks. In terms of the potential broader impact, we do not anticipate any negative ethical and societal impacts of our work while using our method in practice.